\section{The Smith chart}
\label{sec:smith}

\noindent\emph{Status: draft -- the map, the two families of circles and the rotation law are derived; the matching
constructions are stated as chart operations and then verified algebraically in the worked example of
\S\ref{sec:smith-example}. No figure is included: the preamble has no drawing package, so the geometry is given
as algebra plus a table of values.}

A chart is a device for turning an algebraic relation into a visual one, and the Smith chart earns its place in
every microwave course because the relation it turns is the one that appears in every line problem: the reflection
coefficient of a load, seen through a length of line, is the load's reflection coefficient rotated by a phase that
depends on the length. Section~\ref{sec:tline-ripple} derived that rotation, \eqref{eq:tline-gamma-z}; this
chapter takes it seriously. Once $\Gamma$ is drawn as a point inside the unit disc, the two operations that a
line problem needs -- moving along the line, and converting between impedance and admittance -- become a
rotation and a reflection of the page, and matching becomes a construction rather than a search. The cost of
doing it on paper has fallen, since a computer does the same arithmetic in a microsecond, but the understanding
has not: a matching network designed by an optimiser that converges to a solution with a narrow bandwidth is
doing something the chart makes obvious in advance, which is the point of \S\ref{sec:smith-bandwidth}.

The chapter is deliberately self-contained about the geometry and does not appeal to a printed chart. Everything a
chart would let one read off is computed here from the same formula, and the small number of landmarks that are
actually needed are tabulated in Table~\ref{tab:smith-landmarks}.

\subsection{The bilinear map from impedance to reflection coefficient}
\label{sec:smith-map}

Normalise the impedance by the reference, $z=Z/Z_0=r+jx$, so that a matched load is $z=1$ and the two landmarks
of a short and an open are $z=0$ and $z=\infty$. Then \eqref{eq:tline-gamma} reads
\begin{equation}
  \Gamma = \frac{z-1}{z+1},
  \qquad
  z = \frac{1+\Gamma}{1-\Gamma} ,
  \label{eq:smith-map}
\end{equation}
and the inverse in \eqref{eq:smith-map} is obtained by solving for $z$ in one line, which is worth doing once
because it shows that the map is its own kind of symmetry: the same expression, with $\Gamma$ in place of $z$ and
$z$ in place of $\Gamma$, apart from a sign. The map is a M\"obius transformation -- fractional linear, and
composing two of them is again one -- with a single pole at $z=-1$, i.e. at the point $\Gamma=\infty$ which is the
short circuit. It is also the Cayley transform of the right half-plane onto the unit disc, and that identification
is what the next subsection exploits.

Two structural facts follow from \eqref{eq:smith-map} without any computation, and both are what make the chart
usable. First, a M\"obius transformation maps generalised circles to generalised circles: every circle or straight
line in the $z$-plane becomes a circle or straight line in the $\Gamma$-plane, and there is no question of a
family of curves becoming distorted into something unrecognisable. Second, the map is conformal wherever its
derivative does not vanish, and
\begin{equation}
  \frac{\dd\Gamma}{\dd z} = \frac{(z+1)-(z-1)}{(z+1)^2} = \frac{2}{(z+1)^2} \neq 0
  \quad\text{for } z\neq\infty ,
  \label{eq:smith-derivative}
\end{equation}
so it is conformal everywhere except at the open circuit. Conformality means angles are preserved: two curves
that cross at a right angle in the impedance plane still cross at a right angle on the chart. This is why the
constant-resistance and constant-reactance families intersect orthogonally, and it is a property worth stating
because it is what makes reading a chart at a point a reliable operation -- local geometry means what it looks
like it means.

\subsection{Why constant resistance and constant reactance become circles}
\label{sec:smith-circles}

Write $\Gamma=u+jv$ and clear the denominator in \eqref{eq:smith-map}:
\begin{equation}
  (1-\Gamma) z = 1+\Gamma
  \quad\Longrightarrow\quad
  (1-u-jv)(r+jx) = 1+u+jv .
  \label{eq:smith-clear}
\end{equation}
Equating real and imaginary parts gives two real equations,
\begin{equation}
  r - ur + vx = 1+u ,
  \qquad
  x - ux - vr = v ,
  \label{eq:smith-twoeq}
\end{equation}
and eliminating first $x$ and then $r$ yields the two explicit forms
\begin{equation}
  r = \frac{1-u^2-v^2}{(1-u)^2+v^2},
  \qquad
  x = \frac{2v}{(1-u)^2+v^2} ,
  \label{eq:smith-rx}
\end{equation}
which are the equations of the chart written the useful way round: given a point, read off the impedance.
The families of curves are then obtained by fixing $r$ or $x$ in \eqref{eq:smith-rx} and rearranging. For
constant $r$,
\begin{equation}
  r\left[(1-u)^2+v^2\right] = 1-u^2-v^2
  \;\Longrightarrow\;
  \left(u - \frac{r}{1+r}\right)^2 + v^2 = \frac{1}{(1+r)^2} ,
  \label{eq:smith-rcircle}
\end{equation}
a circle with centre $\left(r/(1+r),\,0\right)$ on the real axis and radius $1/(1+r)$. For constant $x$,
\begin{equation}
  x\left[(1-u)^2+v^2\right] = 2v
  \;\Longrightarrow\;
  (u-1)^2 + \left(v - \frac{1}{x}\right)^2 = \frac{1}{x^2} ,
  \label{eq:smith-xcircle}
\end{equation}
a circle with centre $\left(1,\,1/x\right)$ and radius $1/|x|$.

The two families are worth reading as limits. The constant-resistance circles all pass through the point
$(u,v)=(1,0)$, which is the open circuit $\Gamma=+1$: every circle of the family is tangent there, so the right
half of the chart is nested arcs all meeting at one point. As $r\to\infty$ the centre approaches $(1,0)$ and the
radius approaches zero, so the whole real axis of the chart collapses to a point on its rim. The
constant-reactance circles all pass through $(1,0)$ as well, with centres on the vertical line $u=1$; for
$x>0$ they lie in the upper half, for $x<0$ in the lower half, and $x\to0$ degenerates the circle into the real
axis. The orthogonality of the two families is exactly the conformality of \eqref{eq:smith-derivative}: two
circles whose centres lie on perpendicular lines through a common point meet at right angles, and the whole
apparatus of a printed chart is that observation drawn on paper.

\subsection{The passive half-plane is the unit disc}
\label{sec:smith-passive}

The inequality \eqref{eq:tline-gamma-passive} established that a passive load, $\Re Z\ge 0$, gives
$|\Gamma|\le1$. In the $(u,v)$ plane that is the statement that the entire right half of the impedance plane maps
into the closed unit disc, and because the map is a bijection, the correspondence is exact: $r\ge0$ if and only
if $u^2+v^2\le1$. The rim of the disc, $|\Gamma|=1$, is $r=0$, the pure reactances -- the inductors and
capacitors, and the two extremes of short and open at $(-1,0)$ and $(+1,0)$. The real axis of the chart is the
purely resistive line $x=0$, with the matched point $z=1$ at the centre. The upper half of the disc is
$x>0$, the inductive half; the lower half is capacitive. These landmarks, and a few representative points, are
collected in Table~\ref{tab:smith-landmarks}; they are the only ones the constructions of
\S\ref{sec:smith-matching} need, and they replace the figure this report cannot draw.

\begin{table}[H]
  \centering
  \caption{Landmarks of the chart: normalised impedances and the reflection coefficients they map to, computed from \eqref{eq:smith-map}. The pattern along each row is the algebra of \eqref{eq:smith-rcircle} and \eqref{eq:smith-xcircle} evaluated at a fixed value.}
  \label{tab:smith-landmarks}
  \small
  \begin{tabular}{@{}lccc@{}}
    \toprule
    Load & $z=r+jx$ & $\Gamma=u+jv$ & $|\Gamma|$ \\
    \midrule
    matched load        & $1$              & $0$              & $0$ \\
    short circuit       & $0$              & $-1$             & $1$ \\
    open circuit        & $\infty$         & $+1$             & $1$ \\
    pure reactance      & $j1$             & $0.000+j1.000$   & $1$ \\
    pure reactance      & $-j1$            & $0.000-j1.000$   & $1$ \\
    quarter-wave load   & $2$              & $0.333$          & $0.333$ \\
    half the reference  & $0.5$            & $-0.333$         & $0.333$ \\
    the example below   & $0.5+j0.5$       & $-0.2+j0.4$      & $0.4472$ \\
    \bottomrule
  \end{tabular}
\end{table}

The reason the disc is worth drawing at all is that the passivity constraint is a constraint on the picture:
a problem that produces $|\Gamma|>1$ has produced a load that is active, whatever the schematic says. That is the
same inequality that appears as the passivity condition \eqref{eq:sp-passivity} on a scattering matrix, and the
one-port case of it is precisely this rim.

\subsection{Moving along the line is a rotation, one turn per half wavelength}
\label{sec:smith-rotation}

Equation \eqref{eq:tline-gamma-z} says that at a distance $d$ towards the generator,
$\Gamma(d)=\Gamma_L e^{-2j\beta d}$. On the chart this is a rotation of the point about the centre by the angle
$-2\beta d$: a decrease of $2\beta d$ radians in the argument, with the modulus unchanged. The direction matters
and is worth fixing once. Looking towards the generator, $d$ increases, and the rotation is clockwise; looking
towards the load, it is anticlockwise. The modulus is unchanged, so a lossless length of line moves a point
around a circle centred on the origin and never towards it -- which is the geometrical form of the remark at the
end of Section~\ref{sec:tline-zin} that a length of line cannot improve a match by itself.

The period follows from \eqref{eq:tline-ripple-period}. One full turn is $2\pi$, so it requires
$2\beta d=2\pi$, i.e. $d=\pi/\beta=\lambda_g/2$. A complete revolution of the chart is half a guided wavelength,
and a quarter wavelength is a half turn, which is why the quarter-wave transformer \eqref{eq:tline-quarter}
appears on the chart as a rotation by $180^\circ$ taking $z$ to $1/z$. This is the same factor of two as before,
and the reason is the same: the wave goes and comes back. It is worth saying explicitly that the chart's most
familiar scale -- the ``wavelengths towards the generator'' ring printed round the rim -- runs from $0$ to
$0.5$, not from $0$ to $1$, for exactly this reason.

\subsection{VSWR is the distance from the centre}
\label{sec:smith-vswr}

Since moving along a lossless line preserves $|\Gamma|$, and since \eqref{eq:tline-vswr} expresses the standing
wave ratio as a function of $|\Gamma|$ alone, the ratio is constant along the line and corresponds to a circle
centred on the origin. Read \eqref{eq:tline-vswr} the other way and it is transparent:
\begin{equation}
  |\Gamma| = \frac{\mathrm{VSWR}-1}{\mathrm{VSWR}+1} ,
  \label{eq:smith-vswr}
\end{equation}
so the radius of the point is a direct readout of the mismatch. The centre is the perfect match; the rim is total
reflection; a VSWR of $2$ is a radius of $\tfrac13$, and a VSWR of $1.5$ is a radius of $0.2$. This is why the
concentric circles printed on a chart are labelled both in $|\Gamma|$ and in VSWR: they are the same circles. The
circle through the point therefore marks every place along the line where the standing wave has its maxima, and
the crossing of that circle with the real axis towards the load is where the voltage maximum sits -- a fact used
constantly in slotted-line measurements, and the geometrical version of the $\pm1.2$ per cent envelope computed in
Section~\ref{sec:tline-ripple}.

\subsection{Matching, as two constructions}
\label{sec:smith-matching}

Two matching problems account for almost everything the chart is used for, and both are one-line constructions.

The \emph{quarter-wave transformer} uses the half-turn of \S\ref{sec:smith-rotation}. A real load $R_L$ sits on
the real axis at $z=R_L/Z_0$; rotate by $180^\circ$ (a length $\lambda_g/4$) and the point lands at $1/z=Z_0/R_L$,
which by \eqref{eq:tline-quarter} is the impedance seen there. Choose the transformer's characteristic impedance
to be the geometric mean $\sqrt{Z_0 Z_{\mathrm{in}}}$ and both are real, so the match is exact at the design
frequency. The construction is thus: rotate the load by a half turn, read the impedance, and take the square root
of the product. What the construction makes immediately visible is the limit: the half turn is a half turn only at
one frequency, so a quarter-wave transformer is a narrow-band device, and its bandwidth is set by how far the load
is from the centre.

The \emph{shunt stub} uses the same rotation plus one more observation: adding a susceptance at a point means
adding to the imaginary part of the admittance at that point, so the point is moved along a constant-conductance
circle. The construction runs in the admittance plane, which on the chart means reflecting the point through the
centre -- inverting $z$, since $y=1/z$ and $|\Gamma(y)|=|\Gamma(z)|$ with the argument negated (the chart's
``impedance/admittance'' overlay is exactly that reflection). Then: rotate towards the generator until the point
crosses the circle $\Re y=1$, which is the image of the $\Re z=1$ circle under the reflection; at that plane the
input conductance is already matched, and only a susceptance $b$ remains. A shorted stub of length $l<\lambda_g/4$
supplies $y_{\mathrm{stub}}=-j\cot\beta l$, purely susceptive, whose value is read off the chart the same way; set
it to $-b$ and the point lands on the centre. Two properties of the construction are visible and both matter. The
rotation must reach the $\Re y=1$ circle, which it always does on a lossless line, so a solution always exists;
but it may reach it twice, once before and once after the first $\lambda_g/2$, and the two solutions have
different stub lengths and different bandwidths.

\subsection{What the match costs in bandwidth}
\label{sec:smith-bandwidth}

A match designed on the chart is exact at one frequency. Move off it and two things happen at once: the load's
own impedance drifts, and every electrical length in the construction -- the transformer, the stub, the line
between them -- rotates by a different amount than it should. On the chart this is simply the observation that
the point, which was placed on the centre, departs from the centre along some path, and the rate of departure is
what bandwidth means. A construction that reaches the centre from a small radius (a well-matched load) has a
shallow departure and a wide band; one that reaches it from near the rim -- the familiar case of matching a
low-resistance load such as an antenna with a real impedance far from $50\ \Ohm$ -- has a steep departure, because
the rotation angle changes the argument of a vector that was long to begin with.

The quantitative statement is the Bode--Fano limit, which deserves to be quoted for its shape rather than its
derivation: for a load with a fixed reactive part, the integral of $\ln(1/|\Gamma|)$ over frequency is bounded by
a constant that depends only on the load, not on the matching network. The consequence is the one the chart shows
qualitatively: bandwidth and match depth trade off against each other, and a design that achieves a very deep
null over a very narrow band has not done something clever, it has spent the budget in one place. This is the
reason the broadband comparison of Section~\ref{sec:msres-broadband} is worth more than the deepest single-frequency
null in the report: a match at one bin validates one bin, and says nothing about the neighbouring one. When a
matching network is evaluated here, the figure of merit quoted is the band over which the return loss stays below
a stated level, precisely because a single number at the design frequency can be made arbitrarily good.

\subsection{A worked example: $Z_L=25+j25\ \Ohm$ on a $50\ \Ohm$ line}
\label{sec:smith-example}

The example is small enough to check by hand at every step and is carried out completely, because a worked
matching problem is where the conventions of \S\ref{sec:smith-rotation} either become usable or stay abstract.
Take $Z_0=50\ \Ohm$, $Z_L=25+j25\ \Ohm$, and a frequency at which $\lambda_g=100$ mm, corresponding to
$f=1$ GHz on a substrate with $\eps_{\mathrm{eff}}=9$ and $v_p=c/3=10^{8}$ m/s, so that
$\beta=2\pi/\lambda_g=62.832$ rad/m and one millimetre is $3.6^\circ$ of electrical length.

\emph{Step 1: normalise.} $z_L=(25+j25)/50=0.5+j0.5$. The load is in the lower-right quadrant of the chart's
upper half -- more precisely, in the region $0<r<1$, $x>0$, so inside the $r=1$ circle and above the axis.

\emph{Step 2: the reflection coefficient.} From \eqref{eq:smith-map},
\begin{equation}
  \Gamma_L = \frac{(0.5+j0.5)-1}{(0.5+j0.5)+1}
           = \frac{-0.5+j0.5}{1.5+j0.5}
           = \frac{(-0.5+j0.5)(1.5-j0.5)}{1.5^2+0.5^2}
           = \frac{-0.5+j1.0}{2.5} = -0.2+j0.4 .
  \label{eq:smith-ex-gamma}
\end{equation}
Hence $|\Gamma_L|=\sqrt{0.04+0.16}=\sqrt{0.2}=0.4472$, and $\angle\Gamma_L=\arctan(0.4/-0.2)$ in the second
quadrant, $180^\circ-63.435^\circ=116.565^\circ$. The point sits at radius $0.447$ and angle $116.6^\circ$: upper
left, inside the disc, as a passive inductive load must.

\emph{Step 3: the VSWR.} By \eqref{eq:smith-vswr},
\begin{equation}
  \mathrm{VSWR} = \frac{1+0.4472}{1-0.4472} = \frac{1.4472}{0.5528} = 2.6179 ,
  \label{eq:smith-ex-vswr}
\end{equation}
which is the radius of the circle through the point, and the return loss is
$20\log_{10}(0.4472)=-6.99$ dB -- a badly matched load, which is the point of matching it.

\emph{Step 4: the impedance at $\lambda_g/8$.} Here $\beta d=45^\circ$ and the rotation of
\S\ref{sec:smith-rotation} is a quarter turn clockwise:
\begin{equation}
  \Gamma_{\mathrm{in}} = 0.4472\,\angle\,(116.565^\circ-90^\circ) = 0.4472\,\angle\,26.565^\circ
                      = 0.4+j0.2 ,
  \label{eq:smith-ex-gin}
\end{equation}
and inverting \eqref{eq:smith-map},
\begin{equation}
  z_{\mathrm{in}} = \frac{1+0.4+j0.2}{1-0.4-j0.2}
                  = \frac{1.4+j0.2}{0.6-j0.2}
                  = \frac{(1.4+j0.2)(0.6+j0.2)}{0.36+0.04}
                  = \frac{0.80+j0.40}{0.40} = 2.0+j1.0 ,
  \label{eq:smith-ex-zin}
\end{equation}
so $Z_{\mathrm{in}}=100+j50\ \Ohm$. The cross-check through the tangent form \eqref{eq:tline-zin} gives
$\tan45^\circ=1$ and hence $z_{\mathrm{in}}=(0.5+j0.5+j)/(1+j(0.5+j0.5))=(0.5+j1.5)/(0.5+j0.5)=2.0+j1.0$,
agreeing. Two more points on the same circle, for orientation: at $d=\lambda_g/4$ the half turn gives
$z=1/z_L=1/(0.5+j0.5)=1-j1$, i.e. $Z_{\mathrm{in}}=50-j50\ \Ohm$, which is \eqref{eq:tline-quarter} with
$Z_0^2/Z_L=2500/(25+j25)=50-j50$; and at $d=\lambda_g/2$ the full turn returns $Z_{\mathrm{in}}=Z_L$ by
\eqref{eq:tline-half}.

\emph{Step 5: where to put the shunt stub.} Work in admittance. The load admittance is
\begin{equation}
  y_L = \frac{1}{z_L} = \frac{1}{0.5+j0.5} = \frac{0.5-j0.5}{0.5} = 1-j1 ,
  \label{eq:smith-ex-yl}
\end{equation}
whose real part is already unity: the load lies on the $\Re y=1$ circle before any rotation at all, so a stub at
the load plane will do. That is a slightly atypical accident, and the construction is better seen by solving the
general problem, which is where the chart earns its keep. Writing $t=\tan\beta d$, the impedance and admittance
at the plane are
\begin{equation}
  z(d) = \frac{z_L+jt}{1+jz_Lt},
  \qquad
  y(d) = \frac{1}{z(d)} = \frac{1+jz_Lt}{z_L+jt} .
  \label{eq:smith-ex-yd}
\end{equation}
With $z_L=0.5+j0.5$ the numerator is $1+j(0.5+j0.5)t=(1-0.5t)+j0.5t$ and the denominator is $0.5+j(0.5+t)$, so
\begin{equation}
  \Re y(d) = \frac{(1-0.5t)\cdot 0.5 + 0.5t\,(0.5+t)}{0.5 + t + t^2}
           = \frac{0.5+0.5t^2}{0.5+t+t^2} ,
  \label{eq:smith-ex-reyd}
\end{equation}
and imposing $\Re y(d)=1$ gives $0.5+0.5t^2 = 0.5+t+t^2$, i.e.
\begin{equation}
  t^2 + 2t = 0 \quad\Longrightarrow\quad t=0 \ \ \text{or}\ \ t=-2 .
  \label{eq:smith-ex-t}
\end{equation}
Both roots are physical and they are the two crossings of the $\Re y=1$ circle that
\S\ref{sec:smith-matching} promised: $t=0$ is the load plane itself ($d=0$ or $d=\lambda_g/2$), and
$t=-2$ is $\beta d=180^\circ-63.435^\circ=116.565^\circ$, hence $d=(116.565/360)\times100\ \mathrm{mm}=32.38$ mm,
i.e. $d=0.3238\lambda_g$. At the second crossing the admittance is obtained from \eqref{eq:smith-ex-yd} with
$t=-2$:
\begin{equation}
  z(d) = \frac{0.5+j0.5-j2}{1+j(0.5+j0.5)(-2)} = \frac{0.5-j1.5}{2-j1}
       = \frac{(0.5-j1.5)(2+j1)}{5} = \frac{2.5-j2.5}{5} = 0.5-j0.5 ,
  \label{eq:smith-ex-yd-value}
\end{equation}
so $y(d)=1/(0.5-j0.5)=1+j1$, real part unity as required and susceptance $+1$.

Two matching networks therefore exist. At the load plane the residual susceptance is $-1$, so the stub must
supply $+j1$:
\begin{equation}
  -j\cot\beta l = +j1 \quad\Longrightarrow\quad \cot\beta l = -1 \quad\Longrightarrow\quad
  \beta l = 135^\circ , \qquad l = 0.375\lambda_g = 37.5\ \mathrm{mm} .
  \label{eq:smith-ex-stub}
\end{equation}
At the second crossing the residual is $+1$, so the stub must supply $-j1$:
\begin{equation}
  -j\cot\beta l = -j1 \quad\Longrightarrow\quad \cot\beta l = 1 \quad\Longrightarrow\quad
  \beta l = 45^\circ , \qquad l = 0.125\lambda_g = 12.5\ \mathrm{mm} .
  \label{eq:smith-ex-stub2}
\end{equation}
Both give $y_{\mathrm{tot}}=1$ and hence $Z_{\mathrm{in}}=50\ \Ohm$ and $\Gamma=0$ at the design frequency; the
second is preferred on the chart because its stub is shorter, and \S\ref{sec:smith-bandwidth} is the reason the
choice is not merely a matter of taste. The check that closes the loop is a reflection: at the matched plane
$\Gamma$ must be the load point rotated by $-2\beta d=-233.13^\circ$, which is the same as
$+126.87^\circ$, i.e. $\Gamma=0.4472\,\angle\,(116.565^\circ+126.87^\circ)=0.4472\,\angle\,-116.565^\circ
=-0.2-j0.4$. Inverting \eqref{eq:smith-map} gives $z=(0.8-j0.4)/(1.2+j0.4)=0.5-j0.5=1/y(d)$, which is the
reflection of the load point through the centre of the chart -- the impedance/admittance inversion of
\S\ref{sec:smith-matching}, confirmed numerically. Every quantity in the example has now been obtained by two
independent routes and they agree.

\emph{Step 6: what it costs in bandwidth.} The stub supplies exactly $-j1$ at $1$ GHz. At $1.1$ GHz the same
$12.5$ mm is $\beta l=49.5^\circ$ electrically, so $\cot\beta l=0.854$ and the stub supplies only $-j0.854$: a
residual susceptance of $+j0.146$ is left over, and $y=1+j0.146$ corresponds to
$|\Gamma|=0.073$, a return loss of $22.8$ dB. Ten per cent off the design frequency has therefore already cost
the match everything beyond about $-23$ dB, and the trend continues quadratically. This is the numeric content of
\S\ref{sec:smith-bandwidth}: the chart says the same thing in one glance, because the stub's susceptance is read
off a circle and the circle moves fast near its extremes.

That is one line, one match and one tolerance, worked out on the chart. The next section uses all three at once
on a printed line whose reference plane is not where the ports are: the line of Section~\ref{sec:tline}, the
S-parameter and its reference impedance of Section~\ref{sec:sparams} and the chart of this section are the
whole of the vocabulary it needs, and Section~\ref{sec:microstrip} is where they are spent together.
